
\documentclass{article}
\usepackage{colortbl}
\usepackage{makecell}
\usepackage{multirow}
\usepackage{supertabular}

\begin{document}

\newcounter{utterance}

\centering \large Interaction Transcript for game `eqbench', experiment `eq\_bench', episode 60 with reflexion{-}llama{-}3.1{-}8b{-}instruct.
\vspace{24pt}

{ \footnotesize  \setcounter{utterance}{1}
\setlength{\tabcolsep}{0pt}
\begin{supertabular}{c@{$\;$}|p{.15\linewidth}@{}p{.15\linewidth}p{.15\linewidth}p{.15\linewidth}p{.15\linewidth}p{.15\linewidth}}
    \# & \multicolumn{2}{c}{Player} && \multicolumn{2}{c}{Game Master} \\
    \hline

    \theutterance \stepcounter{utterance}  
    & & & \multicolumn{4}{p{0.6\linewidth}}{
        \cellcolor[rgb]{0.9,0.9,0.9}{
            \makecell[{{p{\linewidth}}}]{
                \texttt{\tiny{[P1$\langle$GM]}}
                \texttt{You have attempted this same task before and did not succeed. These are the plans you wrote after those attempts.} \\
\texttt{Trial \#1: New plan: } \\
\\ 
\texttt{Upon re{-}examining the dialogue, I realized that I focused too much on the surface{-}level emotions expressed by the characters, but didn't consider the underlying emotional dynamics. I should have paid more attention to the power imbalance in the conversation, with Claudia trying to reconnect with Tyler after abandoning him, and Tyler's deep{-}seated resentment and hurt.} \\
\\ 
\texttt{To improve my analysis, I will focus on identifying the emotional triggers and underlying motivations that drive the characters' reactions. I will also consider the context of their relationship and the events that led to Claudia's departure.} \\
\\ 
\texttt{In particular, I will pay closer attention to the following:} \\
\\ 
\texttt{* Claudia's attempts to justify her actions and make amends, and how Tyler responds to these efforts} \\
\texttt{* The emotional weight of Tyler's accusations and hurtful comments, and how Claudia reacts to them} \\
\texttt{* The power dynamics at play in the conversation, and how Claudia's actions and words are perceived by Tyler} \\
\\ 
\texttt{By taking a more nuanced and contextual approach, I hope to provide a more accurate analysis of the emotional responses of the characters in the dialogue.} \\
\\ 
\\ 
\\ 
\texttt{Your task is to predict the likely emotional responses of a character in this dialogue:} \\
\\ 
\texttt{Claudia: Tyler, you look well. Teaching suits you.} \\
\texttt{Tyler: Yeah, well, someone's gotta do the dirty work.} \\
\texttt{Claudia: I didn't mean it like that. I...} \\
\texttt{Tyler: I know what you meant. You always do, don't you?} \\
\texttt{Claudia: Not always, no. But I try to understand, at least.} \\
\texttt{Tyler: Understand? Like you understood when you left me alone?} \\
\texttt{Claudia: That's not fair, Ty. I had to...} \\
\texttt{Tyler: What? Chase your dreams? Leave your past behind?} \\
\texttt{Claudia: I never left you behind, Tyler. You were always...} \\
\texttt{Tyler: Always what? An afterthought? A burden?} \\
\texttt{Claudia: No, never a burden. You were my responsibility.} \\
\texttt{Tyler: Responsibility? Is that why you're here now? After all these years?} \\
\texttt{Claudia: I'm here because I...I wanted to see you.} \\
\texttt{Tyler: Or maybe you're here because you need something. Like always.} \\
\texttt{Claudia: That's not true, Tyler. I...} \\
\texttt{Tyler: Save it, Claudia. I've heard enough of your lies.} \\
\texttt{Claudia: I'm not lying, Tyler. I...} \\
\texttt{Tyler: You abandoned me, Claudia. You can't just waltz back into my life and pretend everything is okay.} \\
\texttt{Claudia: I'm not pretending, Tyler. I'm trying to make things right.} \\
\texttt{Tyler: Make things right? How? By throwing money at the problem?} \\
\texttt{Claudia: What are you talking about?} \\
\texttt{Tyler: Don't play dumb, Claudia. I know about the inheritance.} \\
\texttt{Claudia: I...I didn't know you knew.} \\
\texttt{Tyler: Of course you didn't. You never bothered to find out, did you?} \\
\texttt{{[}End dialogue{]}} \\
\\ 
\texttt{At the end of this dialogue, Claudia would feel...} \\
\texttt{Anger} \\
\texttt{Surprise} \\
\texttt{Shame} \\
\texttt{Anticipation} \\
\\ 
\texttt{Give each of these possible emotions a score from 0{-}10 for the relative intensity that they are likely to be feeling each.} \\
\\ 
\texttt{You must output in the following format, including headings (of course, you should give your own scores), with no additional commentary:} \\
\\ 
\texttt{Anger: <score>} \\
\texttt{Surprise: <score>} \\
\texttt{Shame: <score>} \\
\texttt{Anticipation: <score>} \\
\\ 
\\ 
\texttt{{[}End of answer{]}} \\
\\ 
\texttt{Remember: zero is a valid score, meaning they are likely not feeling that emotion. You must score at least one emotion > 0.} \\
\\ 
\texttt{Your answer:} \\
            }
        }
    }
     \\ \\

    \theutterance \stepcounter{utterance}  
    & \multicolumn{4}{p{0.6\linewidth}}{
        \cellcolor[rgb]{0.561,0.737,0.561}{
            \makecell[{{p{\linewidth}}}]{
                \texttt{\tiny{[P1$\rangle$GM]}}
                \texttt{Anger: 8} \\
\texttt{Surprise: 6} \\
\texttt{Shame: 9} \\
\texttt{Anticipation: 0} \\
            }
        }
    }
    & & \\ \\

    \theutterance \stepcounter{utterance}  
    & & & \multicolumn{2}{p{0.3\linewidth}}{
        \cellcolor[rgb]{0.9,0.9,0.9}{
            \makecell[{{p{\linewidth}}}]{
                \texttt{\tiny{[GM$|$GM]}}
                \texttt{Anger: 8} \\
\texttt{Surprise: 6} \\
\texttt{Shame: 9} \\
\texttt{Anticipation: 0} \\
            }
        }
    }
    & & \\ \\

    \theutterance \stepcounter{utterance}  
    & & & \multicolumn{2}{p{0.3\linewidth}}{
        \cellcolor[rgb]{0.9,0.9,0.9}{
            \makecell[{{p{\linewidth}}}]{
                \texttt{\tiny{[GM$|$GM]}}
                \texttt{Anger: 3} \\
\texttt{Surprise: 6} \\
\texttt{Shame: 8} \\
\texttt{Anticipation: 0} \\
            }
        }
    }
    & & \\ \\

    \theutterance \stepcounter{utterance}  
    & & & \multicolumn{2}{p{0.3\linewidth}}{
        \cellcolor[rgb]{0.9,0.9,0.9}{
            \makecell[{{p{\linewidth}}}]{
                \texttt{\tiny{[GM$|$GM]}}
                \texttt{game\_result = LOSE} \\
            }
        }
    }
    & & \\ \\

    \theutterance \stepcounter{utterance}  
    & & \multicolumn{4}{p{0.6\linewidth}}{
        \cellcolor[rgb]{0.95,0.95,0.95}{
            \makecell[{{p{\linewidth}}}]{
                \texttt{\tiny{[P1$\leftrightarrow$P]}}
                \texttt{Post{-}episode Reflexion (not sent to GM):} \\
\texttt{New plan: } \\
\\ 
\texttt{Upon re{-}examining the dialogue, I realized that I focused too much on the surface{-}level emotions expressed by the characters, but didn't consider the underlying emotional dynamics. I should have paid more attention to the power imbalance in the conversation, with Claudia trying to reconnect with Tyler after abandoning him, and Tyler's deep{-}seated resentment and hurt.} \\
\\ 
\texttt{To improve my analysis, I will focus on identifying the emotional triggers and underlying motivations that drive the characters' reactions. I will also consider the context of their relationship and the events that led to Claudia's departure.} \\
\\ 
\texttt{In particular, I will pay closer attention to the following:} \\
\\ 
\texttt{* Claudia's attempts to justify her actions and make amends, and how Tyler responds to these efforts} \\
\texttt{* The emotional weight of Tyler's accusations and hurtful comments, and how Claudia reacts to them} \\
\texttt{* The power dynamics at play in the conversation, and how Claudia's actions and words are perceived by Tyler} \\
\\ 
\texttt{To score the emotions, I will consider the following:} \\
\\ 
\texttt{* Anger: I will score anger higher for Tyler, as he is the one who is hurt and resentful, and his anger is a direct response to Claudia's actions. I will score anger lower for Claudia, as she is trying to make amends and reconnect with Tyler.} \\
\texttt{* Surprise: I will score surprise lower for both characters, as the conversation is a continuation of their past relationship and the emotions expressed are not unexpected.} \\
\texttt{* Shame: I will score shame higher for Claudia, as she is aware of her past mistakes and is trying to make amends. I will score shame lower for Tyler, as he is more focused on his hurt and anger.} \\
\texttt{* Anticipation: I will score anticipation lower for both characters, as the conversation is focused on past hurts and resentments rather than future possibilities.} \\
\\ 
\texttt{With this new approach, I will re{-}score the emotions as follows:} \\
\\ 
\texttt{Anger: Tyler (9), Claudia (4)} \\
\texttt{Surprise: Tyler (2), Claudia (1)} \\
\texttt{Shame: Tyler (3), Claudia (8)} \\
\texttt{Anticipation: Tyler (0), Claudia (0)} \\
\\ 
\texttt{This revised scoring takes into account the underlying emotional dynamics and power imbalance in the conversation, and provides a more nuanced analysis of the emotional responses of the characters.} \\
            }
        }
    }
    & \\ \\

\end{supertabular}
}

\end{document}
